\subsection{Related Work}
\label{sec:relatedwork}
\paragraph{Optimal Taxation.}
In economics, optimal tax theory is the study of the design of a tax system that maximizes a social welfare function subject to a set of economic constraints, while accounting for the fact that individuals respond to taxes and transfers~\citep{mankiw_optimal_2009,diamond_case_2011}.
The core challenge in the design of optimal tax policies is that taxes and transfers can affect
incentives to work, creating a trade-off between equality and productivity~\citep{mankiw_optimal_2009,diamond_case_2011}.
A particular concern is that high income may correlate with high skill, leading  higher skilled workers to choose to work less.


\cite{ramsey1927contribution}'s early  work tied consumption taxes on a good to a representative consumer's elasticity of demand for the good.
The current dominant theoretical framework  arose out of a series of papers by Mirrlees and Diamond~\citep{diamond1971optimal, diamond1971optimal2,mirrlees_optimal_1976}.
These authors consider a utilitarian social planner---aiming to maximize the sum of individual utilities in a society.
\cite{saez_using_2001} builds on the Mirrlees framework to derive optimal non-linear tax rates using models of the elasticity of earnings with respect to tax rates, together with the shape of the income distribution.

Other work has expanded upon the Mirrlees framework to argue for a tax system that tries to achieve a broader distributive justice~\citep{piketty2013theory,piketty2014optimal,saez2016generalized}, or a tax system in which the payments made by an individual merely match the benefits received~\citep{mankiw_optimal_2009,mankiw2010spreading,mankiw2010optimal}.

The Mirrlees model is limited to optimal taxation in a single tax period, without considering dynamics, for example the income histories of individuals in deciding taxes, or events with longer-term effects such as education.
The {\em new dynamic public finance} (NDPF) expands upon these frameworks to consider dynamic economies, capturing additional real world effects, for example, allowing for the coordinated taxation of capital and labor income~\citep{golosov2003optimal,kocherlakota2005zero,albanesi2006dynamic,kocherlakota_new_2010}.

Progress in optimal taxation theory has also come through a growing empirical and experimental literature.
This includes work that seeks to estimate labor supply elasticity to changes in taxation and redistribution~\citep{gruber2002elasticity,chetty2012bounds,goldberg2016kwacha}, and work that seeks to understand the behavioral response of workers to tax policy through the use of cross-sectional data on taxation, labor supply, and individual incomes~\citep{slemrod1996high,goolsbee2000happens,alesina2005work}.
Research in behavioral public finance~\citep{mccaffery2006behavioral,kuziemko2015elastic,alesina2018intergenerational} makes use of experiments and surveys to understand how
people respond to different theories of taxation, redistribution, and public spending.

Our work adopts baselines from  optimal taxation theory, by comparing the performance of the AI Economist with tax policies that arise from the Saez  framework, in this case, making use of estimated labor elasticities in our simulated economies.
\paragraph{Agent-based Modeling.}

Agent-based modeling (ABM) research~\citep{holland1991artificial,bonabeau_agent-based_2002} creates simulations of agents and institutions that interact through prescribed rules.
ABM does not rely on standard equilibrium models. Rather, it allows for dynamic, nonlinear behavior by agents and institutions, and can adopt behavioral rules that are deduced from human experiments~\citep{arthur_designing_1991}.

The idea is to use ABM to enable policy-makers to simulate an artificial economy under different policy scenarios, and quantitatively explore their consequences~\citep{farmer_economy_2009}.
ABM has been applied to study tax compliance~\citep{bloomquist_tax_2011-1,miguel_exploring_2012,subburaj_theory_2018}, and to derive optimal taxation policy~\citep{garrido_agent_2013}, based on heuristics and simple learning methods.
Wider adoption of ABM has proved challenging due to the complexity of realistically modeling human behavior and the economy.

While our motivations are similar to ABM, our framework makes use of deep RL to optimize the behaviors of economic agents with the effect that we study policy design in the presence of rational agent behavior.

\paragraph{Reinforcement Learning.}

Our learning approach relates to {\em multi-agent reinforcement learning} (MARL).
In MARL, agents need to learn together with other learning agents, creating a non-stationary environment~\citep{laurent_world_2011}.
This poses a challenge to the standard approach of learning from exploration~\citep{sutton_reinforcement_2018}, since agents can easily mistake other agents' exploration as environment randomness~\citep{claus_dynamics_1998}.
A particular challenge presented by the AI Economist is that it presents a two-level learning problem, in which the social planner learns a tax policy simultaneously with agents who learn how to optimize their behavior. In effect, agents face a continuously changing reward function. As a consequence, past optimal behavior might not be optimal at later times, which can present a significant learning challenge.

MARL has been effective in learning emergent cooperation in large-scale experiments on complex environments~\citep{bansal_emergent_2017,jaderberg_human-level_2018,OpenAI_dota}.
Previous MARL algorithms have sought to stabilize multi-agent learning by explicitly modeling missing state or policy information~\citep{lowe_multi-agent_2017,tacchetti_relational_2018,shu_m^3rl:_2018}, or assuming some information is shared between agents, including the internal or global state or rewards~\citep{sunehag_value-decomposition_2017,foerster_learning_2017,peysakhovich_prosocial_2017,hughes_inequity_2018,letcher_stable_2018,balduzzi_mechanics_2018}.

In the present paper, we insist on each agent having a policy that only makes use of information that it can individually observe. To make learning efficient, we allow for weight sharing during training. Learned agent behaviors remain distinct, as a result of distinct local states, for example, location in the world, skill, and endowment of resources. This presents a hybrid approach, improving learning efficiency without assuming information  or state sharing between agents.

Optimal taxation can be seen as a form of {\em reward shaping}, which has found a role in preventing undesired social outcomes in multi-agent systems, such as unsustainable resource collection in tragedy-of-the-commons style social dilemmas \citep{leibo_multi-agent_2017}. Reward shaping has also been shown to induce cooperation in spatiotemporal games \citep{mguni_coordinating_2019,hughes_inequity_2018,jaques_social_2018}.
However, these works do not consider the kinds of economic environments we study here,  do not consider the design of tax policies, and  make use of manually-crafted reward shaping.
\paragraph{Machine learning for Economic Design.}
The problem of {\em automated mechanism design} was first formalized by~\cite{ConitzerS02,ConitzerS04a}, and there are polynomial time algorithms for the design of Bayesian incentive-compatible, optimal auctions~\citep{CaiDW12a, CaiDW12b,CaiDW13}.  
  \cite{dutting_optimal_2019} were the first to study the use of deep machine learning for the design of the allocation and payment rules of revenue-optimal auctions. By insisting on incentive-compatible or approximately incentive-compatible designs, their framework can reproduce known optimal designs and also be applied to problems out of reach of current theory. Subsequent work has also adopted neural networks for the design of optimal auctions in settings with budget-constrained bidders~\citep{FengEtAl18},  for the design of  auctions in settings with payment redistribution~\citep{Tacchetti19}, for single-bidder settings~\citep{STZ18}, as well as for problems of social choice~\citep{GolowichEtAl18}.
The use of machine learning for the design of auction mechanisms was earlier pioneered by~\cite{DuettingFJLLP12}, who studied the design of payment rules for a given allocation rule.
Another line of work explores the sample complexity of the problem of learning an optimal auction, typically focusing on simpler settings~\citep{ColeR14,MorgensternR15,BalcanSV16,GonczarowskiW18}. Earlier work studied  the use of machine learning for the design of voting rules~\citep{procaccia09} and for matching and assignment problems~\citep{Narasimhan_ijcai16,Narasimhan_UAI16}.

In the aforementioned settings, the agents do not learn how to behave. Rather, the economic policies (typically auctions) are designed such that  truthful behavior is optimal for an agent.
This avoids the need for two-level learning, where agent behaviors are learned at the same time as an economic policy is learned. An earlier literature did study the co-evolution of  agent behaviors and economic designs, but without making use of reinforcement learning and without studying  tax polices~\citep{DBLP:conf/sigecom/Byde03,DBLP:conf/amec/PhelpsMPS02,DBLP:journals/aamas/PhelpsMP10}. %
Stackelberg equilibria have also been widely studied in other kinds of sequential environments, especially security games. These are two-level problems where the policy of the first-mover (the defender) induces an environment for the second-mover  (the attacker)~\citep{DBLP:conf/atal/PitaJMOPTWPK08,DBLP:books/daglib/0040483}. Recent work has  adopted MARL for the study of security games~\citep{DBLP:conf/aaai/WangSYWSJF19,DBLP:journals/corr/abs-1911-08799}. Two-stage problems also arise in multi-agent  problems where the behavior of some agents is optimized in order to improve the  overall system behavior~\citep{DBLP:conf/nips/DimitrakakisPRT17,DBLP:conf/nips/CarrollSHGSAD19,tylkin20}.
Other work has made use of reinforcement learning to study resource allocation games such as Blotto~\citep{DBLP:conf/icml/BalduzziGB0PJG19}.

Closest in spirit to the present paper, but used for the design of allocation mechanisms  rather than for tax policy (for example matching sellers to buyer queries at Taobao and for internet advertising at Baidu), is the work of~\cite{DBLP:conf/ijcai/Tang17a} and~\cite{shen20}, who make use of RL to improve market design while also allowing for agent behavior to respond to new rules.
\cite{DBLP:conf/aaai/ThompsonNL17} have also advanced the idea of the ``Positronic Economist" (see also~\cite{DBLP:journals/aamas/VorobeychikRW12} and~\cite{DBLP:conf/sigecom/BunzLS18}), which, borrowing from Asimov's positronic brain, describes a system that can be used to represent and then automatically analyze the equilibria that correspond to the rules of economic mechanisms. \cite{parkes_economic_2015} have written, generally, about the role of economic design in economies in which transactions are increasingly mediated through AI systems.